\section{Introduction}
%\vspace{-0.1cm}
\looseness -1 Most reinforcement learning (RL) algorithms are designed to maximize cumulative rewards for a single task at hand. Particularly, model-based RL algorithms, such as~\citep{chua2018pets, kakade2020information, curi2020efficient}, excel in efficiently exploring the dynamical system as they direct the exploration in regions with high rewards. However, due to the directional bias, their underlying learned dynamics model fails to generalize in other areas of the state-action space. While this is sufficient if only one control task is considered, it does not scale to the setting where the system is used to perform several tasks, i.e., under the same dynamics optimized for different reward functions.
As a result, when presented with a new reward function, they often need to relearn a policy from scratch, requiring many interactions with the system, or employ multi-task~\citep{zhang2021survey} or transfer learning~\citep{weiss2016survey} methods. Traditional control approaches such as trajectory optimization~\citep{trajectorPlanningForRobots} and model-predictive control~\citep{mpc} assume knowledge of the system's dynamics. 
They leverage the dynamics model to solve an optimal control problem for each task. 
Moreover, in the presence of an accurate model, important system properties such as stability and sensitivity can also be studied. Hence, knowing an accurate dynamics model bears many practical benefits.
%approaches such as trajectory optimization~\citep{trajectorPlanningForRobots} and model-predictive control~\citep{mpc} are model-based, i.e., they leverage the dynamics model of the system and solve an optimal control problem for each task. 
%However, these methods generally require an accurate dynamics model. This model is often derived using physics' first principles. 
However, in many real-world settings, obtaining a model using just physics' first principles is very challenging. A promising approach is to leverage data for learning the dynamics, i.e., system identification or active learning.
%where a model of the dynamics is learned from data. %Typically, this is also more efficient than standard model-free RL algorithms~\citep{dulac2021challenges}. 
%In many practical settings, it is desirable to first learn the dynamics model of the system, and then use it for control. 
To this end, the key question we investigate in this work is: \emph{how should we interact with the system to learn its dynamics efficiently?} 
    
\looseness -1 
While active learning for regression and classification tasks is well-studied, active learning in RL is much less understood.
    In particular, active learning methods that yield strong theoretical and practical results, generally query data points based on information-theoretic criteria~\citep{krause2008near,settles2009active, balcan2010true,hanneke2014theory, chen15mis}. In the context of dynamical systems, this requires querying arbitrary transitions~\citep{berkenkamp2017safe, mehta2021experimental}.
    However, in most cases, \emph{querying a dynamical system at any state-action pair is unrealistic}. Rather, we can only execute policies on the real system and observe the resulting trajectories. Accordingly, an active learning algorithm for RL needs to suggest policies that are ``informative'' for learning the dynamics. This is challenging since it requires planning with unknown dynamics. %Accordingly,  system identification of nonlinear systems is a widely researched topic, see~\cite{SCHON201139, schoukens} for an overview.  
    %This makes
%system identification of nonlinear systems a very challenging problem, 
%for which theoretical convergence guarantees  only exist for a small class of systems~\citep{simchowitz2018learning, wagenmaker2020active, mania2020active}, see~\cite{SCHON201139, schoukens} for an overview. 
%    While active learning for regression and classification tasks is well-studied~\citep{krause2007nonmyopic, krause2008near}, active learning in RL is much less understood.
 %   Such active learning methods commonly suggest points based on some information theoretic criteria.
  %  Under the assumption that the suggested points are also evaluated, strong theoretical and practical results can be obtained~\citep{balcan2010true,hanneke2014theory}. 
   % However, in RL we {\em cannot query arbitrary transitions}. Rather, we can only execute policies on the real system and observe the resulting trajectories. Accordingly, an active learning algorithm for RL needs to plan and suggest policies with unknown dynamics. %However, in the model-based setting, planning requires knowledge of unknown dynamics.
    
    \paragraph{Contributions}
    In this paper, we introduce a new algorithm, {\em Optimistic Active eXploration (\opacex)}, designed to actively learn nonlinear dynamics within continuous state-action spaces. During each episode, \opacex plans an exploration policy to gather the most information possible about the system. It learns a statistical dynamics model that can quantify its epistemic uncertainty and utilizes this uncertainty for planning. The planned trajectory targets state-action pairs where the model's epistemic uncertainty is high, which naturally encourages exploration. 
    In light of unknown dynamics, \opacex uses an optimistic planner that picks policies that optimistically yield maximal information. We show that this \emph{optimism paradigm} plays a crucial role in studying the theoretical properties of \opacex. Moreover,
 we provide a general convergence analysis for \opacex and prove convergence to the true dynamics for Gaussian process (GP) dynamics models.  Theoretical guarantees for active learning in RL exist for a limited class of systems~\citep{simchowitz2018learning, wagenmaker2020active, mania2020active}, but lack for a more general and practical class of dynamics~\citep{chakraborty2023steering, wagenmaker2023optimal}. 
 We are, to the best of our knowledge, the first to give convergence guarantees for a rich class of nonlinear dynamical systems. 

 We evaluate \opacex on several simulated robotic tasks with state dimensions ranging from two to 58. The empirical results provide validation for our theoretical conclusions, showing that \opacex consistently delivers strong performance across all tested environments.
 Finally, we provide an efficient implementation\footnote{\url{https://github.com/lasgroup/opax}} of \opacex in JAX \citep{jax2018github}.   
  %We propose {\em Optimistic Active eXploration (\opacex)}, an algorithm for actively learning nonlinear dynamics in continuous state-action spaces. At each episode, \opacex executes an exploration policy to gather maximum information about the system. \opacex learns a statistical dynamics model that is capable of quantifying its epistemic uncertainty and leverages the uncertainty to plan an informative trajectory. The informative trajectory visits state-action pairs where the model has high epistemic uncertainty. Thereby, it inherently pursues exploration. In light of unknown dynamics, \opacex uses an optimistic planner which picks policies that optimistically yield maximal information. We show that this \emph{optimism paradigm} is a key component in studying the theoretical properties of \opacex. Moreover,
 %we provide a general convergence analysis for \opacex and prove convergence to the true dynamics for Gaussian process (GP) dynamics models. Theoretical convergence guarantees for system identification only exist for a small class of systems~\citep{simchowitz2018learning, wagenmaker2020active, mania2020active}. 
  % Accordingly,  to the best of our knowledge,  we are the first to give convergence guarantees for a rich class of nonlinear dynamical systems. 
   % We evaluate \opacex on several robotic tasks in simulation with state dimensions from two to 58. We empirically validate our theoretical findings and show that \opacex consistently performs well for all environments.    
